\section{Feasible fractions and the local alternative}
\label{sec:feasible-fractions}

Fix $0<p,\mu,\beta<1$, $w>0$, and $0<x<\min\{p,(1-\mu)^w\}$.
For a row at which both tests fail, conditional Jensen reduces the
moment inequalities to constraints on the two conditional means of $Z$.
We determine their limiting feasible set and maximum weighted mean,
then deduce the local alternative.

\subsection{The two-cell limit}

Suppose a row has $0<q<1$, $m,\pi>0$, and both tests fail.
Set $a=r\mathbb E_\nu(Z\mid Y_i)$ and
$b=r\mathbb E_\nu(Z\mid Y\setminus Y_i)$, taking $b=0$ when the second
cell is empty. Conditional Jensen applied to the convex functions
$(1+z)_+^s$ and $(1-tz)_+^s$, where $t=q/(1-q)$ and $s=\beta r\ge1$,
gives
\begin{equation}\label{eq:finite-cells}
 \begin{aligned}
 m(1+a/r)_+^s+(1-m)(1+b/r)_+^s&\le U(q),\\
 m(1-ta/r)_+^s+(1-m)(1-tb/r)_+^s&\le V(q),\\
 a&>\frac r t\left[1-
       \left(\frac{x^\beta\pi^{1-\beta}V(q)}m\right)^{1/s}\right].
 \end{aligned}
\end{equation}
\leannote{\texttt{graph\_row\_cell\_moments},
\leanfile{S3_TwoCellReduction/L1_ConditionalJensen.lean\#L165}{L1_ConditionalJensen.lean};
the last bound expands \texttt{LF}}
The last inequality follows by taking the positive $s$th root in
the red-cell Jensen bound; it requires only $m>0$.

Consider a sequence satisfying \eqref{eq:finite-cells}, with
$r\to\infty$, $q\to q_0\in(0,\mu]$ and $m,\pi\to d\in[p,1]$.
The lower threshold for $a$ tends to
\[
 A_d(q_0)=\frac{1-q_0}{q_0}\log\frac{d(1-q_0)^w}{x}.
\]
\leannote{\texttt{tendsto\_red\_threshold},
\leanfile{S3_TwoCellReduction/L4_DensityGap.lean\#L150}{L4_DensityGap.lean}}
Indeed, for a positive quantity $K$ tending to $K_0>0$,
$r(1-K^{1/(\beta r)})\to-\log K_0/\beta$.
The first blue-cell moment bounds $a$ above by
$r[(U(q)/m)^{1/(\beta r)}-1]$. Comparing with the convergent lower
threshold gives $d e^{\beta A_d(q_0)}\le U(q_0)$ without assuming
convergence of $a$. Since $x<(1-\mu)^w$, the density-gap argument
of~\cite[Proposition~3.5, Step~6]{ZBS} implies $d<1$.
See \texttt{first\_cell\_limit} and \texttt{density\_log\_gap} in
\leanfile{S3_TwoCellReduction/L4_DensityGap.lean\#L49}{L4_DensityGap.lean}.
Both cell masses are therefore bounded away from zero eventually.
The moment inequalities give
\[
 -\frac r t\left[\left(\frac{V(q)}m\right)^{1/(\beta r)}-1\right]
 \le a\le
 r\left[\left(\frac{U(q)}m\right)^{1/(\beta r)}-1\right].
\]
\leannote{\texttt{scaled\_mean\_bounds},
\leanfile{S3_TwoCellReduction/L2_FiniteBounds.lean\#L32}{L2_FiniteBounds.lean}}
The same bounds hold for $b$ with $m$ replaced by $1-m$.
All four endpoints remain bounded. We may therefore extract a common
subsequence along which $a\to A$ and $b\to C$.
Passing to the exponential limits in \eqref{eq:finite-cells} yields,
writing $q=q_0$ and $t=q/(1-q)$,
\begin{equation}\label{eq:feasible-pair}
 \begin{aligned}
 A&\ge A_d(q),\\
 d e^{\beta A}+(1-d)e^{\beta C}&\le U(q),\\
 d e^{-\beta tA}+(1-d)e^{-\beta tC}&\le V(q).
 \end{aligned}
\end{equation}
\leannote{\texttt{exists\_feasible\_subsequence\_of\_row\_moments},
\leanfile{S3_TwoCellReduction/L5_LimitingConstraints.lean\#L135}{L5_LimitingConstraints.lean}}
Call such a pair \emph{feasible at density $d$ and fraction $q$}.
For $0<d\le1$, let $F_d$ be the set of $q\in(0,\mu]$ admitting a
feasible pair.

\subsection{Maximizing the weighted mean}

Write $B_q=\beta^{-1}\log U(q)=w\log(\mu/q)$.
The complementary constraint changes the set $F_p$, but leaves the
blue-moment optimizer unchanged for every $q\in F_p$.

\begin{lemma}\label{lem:optimizer}
For $q\in F_p$, the maximum of $pA+(1-p)C$ is
\begin{equation}\label{eq:optimizer}
 M_p(q)=
 \begin{cases}
 B_q,& A_p(q)\le B_q,\\[1mm]
 pA_p(q)+\dfrac{1-p}{\beta}
 \log\dfrac{U(q)-p e^{\beta A_p(q)}}{1-p},& A_p(q)>B_q.
 \end{cases}
\end{equation}
\leannote{\texttt{optimal\_envelope\_attained},
\leanfile{S4_ScalarOptimization/L6_OptimalEnvelope.lean\#L76}{L6_OptimalEnvelope.lean}}
Moreover, if $p\le d\le1$ and $(A,C)$ is feasible at density $d$,
then $q\in F_p$ and $dA+(1-d)C\le M_p(q)$.
\end{lemma}
\begin{proof}
Jensen gives $e^{\beta(pA+(1-p)C)}\le U(q)$, so the mean is at most
$B_q$. If $A_p(q)\le B_q$, the pair $(B_q,B_q)$ attains this bound.
It also satisfies the complementary constraint, since $B_q\ge0$ and
$e^{-\beta tB_q}\le1\le V(q)$.

Suppose instead that $A_p(q)>B_q$, and start with any feasible pair.
Since $(1-p)e^{\beta C}>0$, we have $p e^{\beta A_p(q)}<U(q)$.
Increase $e^{\beta C}$ until the blue moment equals $U(q)$.
This increases the logarithmic objective and decreases the
complementary moment. Now contract $e^{\beta A}$ and $e^{\beta C}$ linearly
toward their weighted mean $U(q)$ until the first reaches
$e^{\beta A_p(q)}>U(q)$. Convexity of $z\mapsto z^{-t}$ again decreases
the complementary moment; concavity of $\log z$ increases the objective.
The resulting first coordinate is $A_p(q)$ and the second is determined by
the blue equality. This proves the second branch and its attainment.
If this canonical pair fails the complementary bound, no feasible pair exists.

To remove the unknown limiting density $d$, keep $A$ and define $C'$ by
\[
 e^{\beta C'}=\frac{(d-p)e^{\beta A}+(1-d)e^{\beta C}}{1-p}.
\]
\leannote{\texttt{feasible\_mass\_transfer},
\leanfile{S4_ScalarOptimization/L3_DensityTransfer.lean\#L25}{L3_DensityTransfer.lean}}
This moves mass $d-p$ from the first cell into the second and preserves
the blue moment. Jensen for $z^{-t}$ shows that the complementary moment
does not increase. Jensen for $\log z$ gives
$dA+(1-d)C\le pA+(1-p)C'$. Finally, $A_d(q)\ge A_p(q)$, so the new pair
is feasible at density $p$ and the preceding optimization applies.
\end{proof}

\subsection{A convex test for feasibility}

Put $\theta=p(x/p)^\beta\in(0,p)$ and define
\begin{equation}\label{eq:entropy-test}
 \begin{aligned}
 Z_\theta(q)&=\theta(p/\theta)^{1/q}
       +(1-\theta)((1-p)/(1-\theta))^{1/q},\\
 \mathcal F(q)&=q\log Z_\theta(q)
       +\beta w\bigl[(1-q)\log(1-q)+q\log(q/\mu)\bigr].
 \end{aligned}
\end{equation}
\leannote{\texttt{entropyGap},
\leanfile{S4_ScalarOptimization/L7_EntropyFeasibility.lean\#L80}{L7_EntropyFeasibility.lean}}

\begin{lemma}\label{lem:feasibility-cutoff}
For $0<q\le\mu$, we have $q\in F_p$ if and only if $\mathcal F(q)\le0$.
The function $\mathcal F$ is convex on $(0,1)$. Consequently, if
$0<a<\mu$ and $\mathcal F(a)>0>\mathcal F(\mu)$, then
$\mu\in F_p\subset(a,\mu]$.
\end{lemma}
\begin{proof}
Set
$z_1=e^{-\beta tA}/V(q)$ and $z_2=e^{-\beta tC}/V(q)$.
Since $p e^{-\beta tA_p(q)}/V(q)=\theta$, their constraints become
$z_1\le\theta/p<1$ and $pz_1+(1-p)z_2\le1$.
Increase $z_2$ until equality holds, then contract both coordinates toward
the mean $1$ until $z_1=\theta/p$. Because $z^{-1/t}$ is decreasing
and convex, these operations minimize the blue moment among all pairs
satisfying the complementary constraint and the threshold.
The minimizing coordinates are $\theta/p$ and $(1-\theta)/(1-p)$.
Their blue moment is $V(q)^{-1/t}Z_\theta(q)$.
Taking logarithms of its comparison with $U(q)$, and multiplying by
$q>0$, gives exactly $\mathcal F(q)\le0$.

For convexity, $\log Z_\theta(q)$ is a convex log-sum-exp evaluated at
$1/q$; its positive perspective $q\log Z_\theta(q)$ is convex.
The remaining term is a positive multiple of
$q\log q+(1-q)\log(1-q)$, plus an affine term, and is also convex.
If a feasible $q\le a$ existed, convexity between $q$ and $\mu$ would
give $\mathcal F(a)\le0$, a contradiction. The sign at $\mu$ proves
nonemptiness.
\end{proof}
\leannote{the feasibility equivalence in
\leanfile{S4_ScalarOptimization/L7_EntropyFeasibility.lean\#L242}{L7_EntropyFeasibility.lean}
and the cutoff in \leanfile{S4_ScalarOptimization/L8_ConvexDegreeCutoff.lean\#L172}{L8_ConvexDegreeCutoff.lean}}

\subsection{Uniform approximation and the local alternative}

The next lemma gives a uniform finite approximation to the limiting bound.

\begin{lemma}\label{lem:finite-approximation}
For every $\delta,\varepsilon>0$, there exist
$0<\eta<\min\{p/2,1-\mu\}$ and $R>1/\beta$ such
that, if $r\ge R$, the bounds \eqref{eq:finite-cells} hold,
$\delta\le q\le\mu+\eta$, $m,\pi\in[p-\eta,1]$, and
$|m-\pi|\le\eta$, then some $q'\in F_p$ satisfies
\[
 |q'-q|\le\varepsilon,\qquad
 q\bigl(ma+(1-m)b\bigr)\le q'M_p(q')+\varepsilon.
\]
\leannote{\texttt{finite\_row\_approximation},
\leanfile{S5_LocalAlternative/L1_FiniteApproximation.lean\#L25}{L1_FiniteApproximation.lean}}
\end{lemma}
\begin{proof}
Otherwise choose counterexamples with $r\to\infty$ and $\eta\to0$.
After passing to a subsequence, $q\to q_0\in[\delta,\mu]$ and
$m,\pi\to d\in[p,1]$. The density gap and moment bounds above give
$d<1$ and convergent scaled means. Their limiting pair is feasible
at $d$, so Lemma~\ref{lem:optimizer} gives $q_0\in F_p$ and bounds the
limiting objective by $q_0M_p(q_0)$. Taking $q'=q_0$ contradicts the
assumed failure of the conclusion for all sufficiently late terms.
\end{proof}

\begin{proposition}[Local alternative]\label{prop:local-alternative}
Suppose $F_p\ne\varnothing$ and $E:(0,1]\to\mathbb R$ is continuous,
nondecreasing and integrable, with
\begin{equation}\label{eq:negative-envelope}
 qM_p(q)+wq(1+\log u)\le E(u)
 \quad(q\in F_p,\ 0<u\le1),\qquad
 \int_0^1E(u)\,du<0.
\end{equation}
\leannote{\texttt{ComplementEnvelope},
\leanfile{S4_ScalarOptimization/L1_FeasiblePairs.lean\#L31}{L1_FeasiblePairs.lean}}
There exist $\eta,\tau,\zeta>0$ and $R>\max\{w,1/\beta\}$ such that,
for every $r\ge R$, every maximal positive-potential configuration with
$|c-p|\le\eta$ has the following property. If $X_0\subseteq X$,
$|X\setminus X_0|\le\zeta|X|$ and $q_i\le\mu+\tau$ for $i\in X_0$,
then either a blue test succeeds at some $i\in X$ with $q_i>0$, or a
shifted-red test succeeds at some $i\in X_0$.
\end{proposition}
\begin{proof}
Suppose all the indicated tests fail. The concentration and small-degree
estimates of~\cite[Proposition~3.5, Steps~2 and~5]{ZBS}, using $x<p$,
give a fixed $\delta>0$ with the following property. For any sufficiently
small $\varepsilon,\varphi>0$, choose the configuration tolerances small
and $r$ large; all but $\varphi|X|$ rows then lie in $X_0$ and satisfy
Lemma~\ref{lem:finite-approximation}'s hypotheses. Each row outside this
exceptional set therefore has $q'_i\in F_p$ with $|q'_i-q_i|\le\varepsilon$
and $q_i(m_i a_i+(1-m_i)b_i)\le q'_iM_p(q'_i)+\varepsilon$.

Put $N=|X|$ and $P_i=q_i\mathbb E_\nu Z_i$, taking $P_i=0$ when
$q_i=0$. The identity $rP_i=q_i(m_i a_i+(1-m_i)b_i)$ identifies
the finite objective with the row contribution.
Order the rows by increasing $q_i$, and let $\Delta L_i$ be the
increment of $L(u)=u\log u$, with $L(0)=0$, on $[i/N,(i+1)/N]$.
Stationarity and ordered drift~\cite[Proposition~3.5, Steps~3--4]{ZBS}
give the lower bound
\[
 T:=\sum_{i=0}^{N-1}\left(\frac{rP_i}{N}+wq_i\Delta L_i\right)
 \ge-\frac{4w^2}{r}\qquad(r\ge2w).
\]
\leannote{\texttt{ordered\_envelope\_contradiction},
\leanfile{S5_LocalAlternative/L2_OrderedEnvelope.lean\#L61}{L2_OrderedEnvelope.lean},
using the inherited \texttt{lower\_sum}}
For each nonexceptional row $i$, integrate \eqref{eq:negative-envelope}
with $q'_i$ over $[i/N,(i+1)/N]$. Since $|L'(u)|\le1-\log u$, the total
approximation error is at most $\varepsilon(1+2w)$, including the
interval with left endpoint zero. Let $\mathcal B$ be the union of these
intervals for exceptional rows, of total length at most $\varphi$. Comparing there with
any fixed $q_*\in F_p$ and using the blue-failure bound $rP_i\le w$
gives
\[
 T\le\int_0^1E(u)\,du+\varepsilon(1+2w)
 +(w+|q_*M_p(q_*)|)\varphi
 +w\int_{\mathcal B}(1-\log u)\,du.
\]
\leannote{the estimates for nonexceptional and exceptional rows in
\leanfile{S5_LocalAlternative/L2_OrderedEnvelope.lean\#L122}{L2_OrderedEnvelope.lean};
the integrals are its increments of \texttt{JF}}
The last two terms tend to zero with $\varphi$, uniformly in $N$,
by integrability of $1-\log u$. Choose $\varepsilon,\varphi$ small,
then the configuration tolerances small and $r$ large. The resulting
fixed negative upper bound contradicts $T\ge-4w^2/r$ for sufficiently
large $r$.
\end{proof}
\leannote{the local alternative,
\leanfile{S5_LocalAlternative/L3_GraphAlternative.lean\#L283}{L3_GraphAlternative.lean}}
